\section{Giới thiệu bài toán và mục tiêu}

Ransomware (mã độc tống tiền) mã hoá dữ liệu của nạn nhân và đòi tiền chuộc, thường được trả bằng Bitcoin. Vì mọi giao dịch Bitcoin
đều công khai trên chuỗi khối, ta có thể quan sát hành vi của các địa chỉ nhận tiền và đặt câu hỏi: \emph{chỉ từ các đặc trưng giao dịch
của một địa chỉ trong một ngày, có thể nhận ra địa chỉ đó là địa chỉ bình thường hay thuộc một họ ransomware cụ thể không?}

{\sloppy\paragraph{Bài toán.} Phân loại đa lớp (7 lớp): \code{white} (bình thường), 5 họ ransomware phổ biến nhất
(\code{paduaCryptoWall}, \code{montrealCryptoLocker}, \code{princetonCerber}, \code{princetonLocky}, \code{montrealCryptXXX})
và \code{otherRansom} (gộp 23 họ hiếm còn lại). Cách định nghĩa nhãn, đặc trưng và cách chia dữ liệu được giữ nguyên như bài giữa kỳ
(Bài~1) để các kết quả so sánh được với nhau.\par}

\paragraph{Mục tiêu đồ án.}
\begin{enumerate}
\item Xây dựng, huấn luyện và đánh giá mạng nơ-ron hồi tiếp (LSTM, GRU) và mạng tích chập 1 chiều (CNN Conv1D) cho bài toán trên.
\item Vì dữ liệu ở dạng bảng, đề xuất và so sánh \textbf{hai cách biểu diễn chuỗi}:
  (A) coi 11 đặc trưng của một dòng là một chuỗi;
  (B) dựng \emph{chuỗi lịch sử thật} của mỗi địa chỉ qua các ngày.
\item Thực nghiệm có hệ thống ($\approx$50 cấu hình, các cấu hình chính chạy 3 seed), so sánh với các mô hình nền (Logistic Regression, KNN, Random Forest,
  HistGradientBoosting, LightGBM, XGBoost, MLP), phân tích overfitting/underfitting và các lỗi dự đoán.
\item Đánh giá trung thực: chỉ ra khi nào RNN/CNN \emph{không} phải lựa chọn tốt nhất và vì sao.
\end{enumerate}
